% LNCS systems paper on the ONNX-EdgeIDS software. Lives in the software's own
% repository (github.com/vuthainguyen1602/onnx-edge-ids) under lncs/; the short
% Software Impacts article about the same software is a separate manuscript
% kept outside the repository.
% Compile: ./compile.sh
%
% STATUS: draft. Target venue: DIMVA 2027 (cycle 2 ~Feb 2027 expected);
% backup: ACNS 2027 cycle 2 (Jan 21, 2027). See ../README.md.
% Single author; the co-authors of the earlier companion paper are not on
% this manuscript.
%
% RESULT PROVENANCE: the cold-start, throughput, latency, memory, power, and
% parity numbers in this draft are the same validation-run figures recorded
% in onnx-edge-ids/submission_checklist.md (single replay run, Jetson #2 on
% Wi-Fi). Re-verify against that file (or a fresh run) before submission; the
% End-to-end latency was measured over Wi-Fi on 2026-10-11 (exp7); the
% remaining \ph{...} marks figures still to be repeated and must not be invented. The traversal and early-exit timings were
% re-measured on Jetson #2 on 2026-10-11 (../code/results_jetson/). Table tab:boundary, Fig. fig:lockstep, Tables tab:earlyexit and tab:generalize
% are regenerated by ../code/experiments/ (exp2/exp2b, exp1, exp3, exp5/exp5b)
% and quoted from ../code/results/*.json.
%
\documentclass[runningheads]{llncs}

\usepackage{fontspec}
\usepackage{polyglossia}
\setdefaultlanguage{english}
\setotherlanguage{vietnamese}
% Shared XeLaTeX font setup for Vietnamese manuscripts.
% TeX Gyre Termes is a Times-compatible font that ships with TeX Live / MacTeX,
% covers Vietnamese, and needs no MS-Office fonts -> builds on any machine.
% If TeX Gyre Termes is missing (e.g. BasicTeX without the tex-gyre package),
% fall back to macOS' built-in Times New Roman instead of silently dropping to
% nullfont (which produces blank pages).
\IfFontExistsTF{TeX Gyre Termes}
  {\setmainfont{TeX Gyre Termes}[Ligatures=TeX]}
  {\setmainfont{Times New Roman}[Ligatures=TeX]}

\usepackage{graphicx}
\usepackage{booktabs}
\usepackage{amsmath}
\usepackage{xcolor}
\usepackage{url}
% Okabe-Ito colorblind-safe accents for TikZ diagrams
\definecolor{accBlue}{HTML}{0072B2}
\definecolor{accGreen}{HTML}{009E73}
\definecolor{accOrange}{HTML}{D55E00}
\usepackage{tikz}
\usetikzlibrary{positioning, arrows.meta, shapes.geometric, fit, backgrounds, calc}
\usepackage{pgfplots}
\pgfplotsset{compat=1.18}
\usepgfplotslibrary{groupplots}

\newcommand{\ph}[1]{{\color{gray}\textit{[#1]}}}
\newcommand{\tabnote}[1]{%
  \vspace{4pt}\par\noindent
  \begin{minipage}{\textwidth}%
    \footnotesize\parindent=0pt\leftskip=0pt\rightskip=0pt\parfillskip=0pt plus 1fil\relax
    #1%
  \end{minipage}}
\renewcommand{\topfraction}{0.92}
\renewcommand{\bottomfraction}{0.6}
\renewcommand{\textfraction}{0.06}
\renewcommand{\floatpagefraction}{0.85}
\setlength{\floatsep}{8pt plus 2pt minus 2pt}
\setlength{\textfloatsep}{10pt plus 4pt minus 3pt}
\setlength{\emergencystretch}{1.5em}
\providecommand{\babelshorthandoff}[1]{}

\graphicspath{{../results/benchmarks/}{../figures/}}

\begin{document}

\renewcommand{\figurename}{Fig.}

\title{Two Runtimes, One Verdict: Boundary-Directed Parity Testing for Tree-Ensemble Intrusion Detection at the Edge}

\titlerunning{Boundary-Directed Parity for Edge Intrusion Detection}

\author{Thai Nguyen Vu}

\authorrunning{T. N. Vu}

\institute{Campus in Ho Chi Minh City, University of Transport and Communications, Ho Chi Minh City, Vietnam\\
\email{thaivn\_ph@utc.edu.vn}}

\maketitle

\begin{abstract}
An edge intrusion detector that ships a primary and a fallback runtime promises the same verdict from both. We test it on a two-stage Kafka pipeline on Jetson Orin Nano Super boards whose 200-tree random forest is served by ONNX Runtime or by a NumPy engine exported from the same ONNX graph. Replay is a weak test: tree runtimes can only disagree when a feature value lands on a split threshold, and recorded flows almost never do. We derive one probe per split node, on its threshold and one float32 step either side. A float64 NumPy traversal flips 328 labels against ONNX Runtime on the probes but 15 on the whole CICIDS2017 test split, so a 2{,}000-row replay gate passes it 93.5\% of the time. The deployed float32 traversal flips none. Flips reach the label because the trainer's 32-bin split search leaves the forest's splits sharing only 815 distinct thresholds. On four further forests from scikit-learn, LightGBM and XGBoost the same probes expose traversal and converter divergences 3 to $169\times$ more densely than replay, while only replay catches two ONNX Runtime releases that relabel 85\% of flows for a stock RandomForest export. A block-lockstep traversal ($80\times$ faster at batch size 1 on the board) and an exact early-exit rule (half the trees, no label changed) make the fallback usable in streaming. Frozen on CSE-CIC-IDS2018 traffic the detector collapses, both serving guarantees hold on every shifted flow, and the share of flows early exit cannot settle, calibrated as a label-free drift alarm, rises from 4.8\% to 28.2\%, a label-free drift signal.

\keywords{Intrusion detection \and Edge computing \and ONNX Runtime \and Differential testing \and Random forest \and Apache Kafka}
\end{abstract}

\section{Introduction}

Streaming ML-based intrusion detection increasingly runs at the network edge to cut bandwidth and response time~\cite{shi2016edge,hamdan2020edge}. A recurring failure mode in such deployments lies in the serving stack around the model rather than in the model itself: a framework convenient for offline training is reused for online serving because it is operationally simpler to run one stack than two, and that convenience carries a memory and latency cost that a constrained ARM board cannot always absorb. Earlier deployment work on this project's two-stage anomaly-gate-then-classifier pipeline exposed exactly this limitation: the classifier was accurate, but its serving path was heavier than the 8\,GB Jetson Orin Nano Super boards it ran on could comfortably sustain.

The response to that limitation was \textsc{ONNX-EdgeIDS}\footnote{\url{https://github.com/vuthainguyen1602/onnx-edge-ids}}, a serving toolkit organized around one portable artifact, an ONNX \texttt{TreeEnsembleClassifier} graph, with two interchangeable runtimes behind it. The primary runtime is ONNX Runtime. A second, dependency-minimal runtime re-implements the same forest in pure NumPy for boards where no matching \texttt{onnxruntime} wheel exists. The NumPy exporter does not depend on the framework that produced the ONNX graph: it parses the graph's standard tree-ensemble operator attributes (split features, thresholds, per-leaf class weights), not a framework-specific internal representation, so the same exporter works on an ONNX tree ensemble regardless of its origin. A companion validation script treats agreement between the two runtimes as a release gate instead of a one-off sanity check, failing the build if label agreement or confidence drift falls outside tolerance.

We integrate both runtimes into a two-stage Kafka pipeline (an anomaly gate on one Jetson board forwarding suspicious flows to a classifier on a second) and then ask three questions that the integration itself raises. \emph{Do the two runtimes really agree?} Replaying recorded flows says yes, but recorded flows almost never sit on a split threshold, which is the only place two tree runtimes can disagree. We probe every threshold directly instead (Sect.~\ref{sec:boundary}). \emph{Why is the NumPy fallback slow at the small batches a streaming node serves?} Because it pays Python dispatch once per tree per level. Advancing blocks of trees in lockstep removes that (Sect.~\ref{sec:lockstep}). \emph{How much of the forest does forwarded traffic need?} Barely more than half, under a rule that cannot change a label (Sect.~\ref{sec:earlyexit}). A fourth question follows from the first: whether its answer depends on this forest at all, which Sect.~\ref{sec:generalize} tests on four others.

\textbf{Contributions.}
(1) \emph{Boundary-directed parity testing} for tree-ensemble serving artifacts: one probe per split node, on the threshold and one float32 step either side, with a path-satisfying construction that reaches 99.9\% of nodes, where the same boundary values without routing exercise 27\% and find $2.5\times$ fewer flips. On the deployed forest it exposes a float64/float32 traversal divergence $34\times$ more densely than replaying the whole labelled test split, and shows that a 2{,}000-row replay gate would miss it 93.5\% of the time. We also explain why such divergences flip \emph{labels} rather than single votes: the trainer's binned split search leaves 95{,}741 splits sharing 815 distinct thresholds, so arithmetic disagreement at one value is correlated across the ensemble, and the same fact makes an exhaustive boundary gate cheap. On four further forests from three other trainers the probes generalize, and a replay-visible ONNX Runtime version regression shows what they are not for;
(2) A block-lockstep NumPy traversal, measured on the board on real flows against the per-tree loop it replaces, across batch sizes 1 to 2{,}000;
(3) An exact early-exit rule for the binary forest and its measured cost on gate-forwarded traffic;
(4) A shift stress test that separates what drift breaks (the detector and the gate) from what it cannot (runtime parity and early-exit exactness), and calibrates the early-exit undecided rate as a label-free drift alarm: on windows of 10{,}000 flows it raises no false alarm on held-out in-distribution windows and fires on every shifted one.
The framework-agnostic ONNX-to-NumPy exporter, the on-board measurements of both runtimes, and the experiment code, artifacts and result files behind every number in this paper are public\footnote{\url{https://github.com/vuthainguyen1602/onnx-edge-ids} (directory \texttt{lncs/code/}); software archive of the measured version, v0.3.0: \url{https://doi.org/10.5281/zenodo.23299232}.}.

\section{Related Work}

\textbf{Edge intrusion detection.} Single-board deployments of ML-based IDS are now common, from early Raspberry Pi systems~\cite{sforzin2016rpids} to recent ensemble and metaheuristic-optimized detectors on IoT-edge and industrial-control hardware~\cite{baalaji2025ensemble,asghar2026hedid,zhang2026edgeics} and autoencoder- or generative-model-based detection on Jetson boards~\cite{hasan2025iiot,diazgorrin2026jetson}. These studies optimize the detector. The present paper instead treats the \emph{serving stack} as the object of study, on the observation (from earlier deployment work on this system) that model accuracy and serving cost are separate axes that a single evaluation can conflate.

\textbf{Serving classical ML on constrained hardware.} ONNX\footnote{\url{https://onnx.ai/}} standardizes a small operator set, including \texttt{TreeEnsembleClassifier} for decision-tree ensembles, so a model trained in one framework can be served by any ONNX-compatible runtime. ONNX Runtime\footnote{\url{https://onnxruntime.ai/}} is the reference CPU/GPU implementation and is what we target as the primary backend. Nakandala et al.~\cite{nakandala2020hummingbird} make a related argument for tree ensembles specifically: compiling scikit-learn-style pipelines into tensor operations lets a GPU tensor runtime, rather than a tree-optimized one, serve them, with large reported speedups over native execution on batched workloads. Our NumPy fallback takes the opposite trade-off deliberately: no compilation step and no GPU dependency, at the cost of the batch-size sensitivity this trades in. Koschel et al.~\cite{koschel2023arm} show that this cost is not fundamental to ARM hardware: porting the QuickScorer tree-traversal algorithm~\cite{lucchese2015quickscorer} from Intel AVX to ARM NEON yields up to $9.4\times$ speedups for ensemble inference on IoT-class ARM devices, with fixed-point quantization improving throughput further at a small accuracy cost. We do not port QuickScorer's bitvector algorithm or write NEON intrinsics. We take its underlying lesson, evaluate many trees in lockstep instead of one at a time, apply it at the NumPy level to the NumPy engine's traversal loop, and measure the result (Sect.~\ref{sec:lockstep}).

\textbf{Two-stage detection, ensemble early exit, and runtime equivalence.} Cheap-filter-then-forest cascades are now routine in intrusion detection, including on embedded targets~\cite{ferreira2026canbus}. Our pipeline has that shape and we claim nothing new for it. Stopping an ensemble before every member has voted is also old: Hern\'andez-Lobato et al.~\cite{hernandezlobato2009pruning} halt once the remaining votes are statistically unlikely to change the outcome, and Wang et al.~\cite{wang2018qwyc} learn an evaluation order and per-step thresholds jointly. Both trade a small, tunable disagreement for speed. An intrusion detector with two runtimes cannot take that trade without making the runtimes disagree with each other, so Sect.~\ref{sec:earlyexit} uses the exact counterpart, stopping only when the unevaluated trees cannot flip the label, and measures what the exactness costs. On equivalence, differential testing is well established for deep learning. CRADLE~\cite{pham2019cradle} runs one model on several backends and localizes the library bug behind each inconsistency. Wen et al.~\cite{wen2025mindthegap} do the same across five accelerator platforms, and Chen~\cite{chen2026integeralibi} traces bitwise divergence between two INT8 kernels serving one language model to the stage where it arises. Cross-framework fuzzing shows that boundary and non-finite inputs expose inconsistent comparison semantics~\cite{duan2026xamt}, and Jajal et al.~\cite{jajal2024onnxconverters} catalogue the failures of ONNX model converters, almost all of them on deep networks. For tree ensembles, Naik et al.~\cite{naik2026sensitivity} analyze symbolically how far the output can move under small input perturbations, and verifiers such as Veritas~\cite{devos2021veritas} prove properties of one ensemble over hyperrectangles of the input space. Both reason over the reals. The disagreements in this paper sit in the last bit of a float32 at a threshold, where the two runtimes' arithmetic differs. The precision sensitivity of tree thresholds under float32 export is documented by the converter projects themselves: the sklearn-onnx tutorial on switching from double to float shows that a tree, being a discontinuous function of its input, can change its prediction under a rounding-sized perturbation~\cite{sklearnonnx_float}. We have not found it applied where it matters operationally: as a release gate for two \emph{serving} runtimes of one deployed detector, with probes derived from the model's own split table rather than from recorded or random inputs.

\textbf{Evaluation protocol.} Arp et al.~\cite{arp2022dosdonts} and Lanvin et al.~\cite{lanvin2023errors} document how easily an offline security-ML evaluation is inflated by leakage or dataset artifacts. The classifier served by this toolkit was selected under a leakage-aware, deduplicated CICIDS2017~\cite{sharafaldin2018toward} protocol with SHAP-based~\cite{lundberg2017unified} feature selection in earlier work on this system, out of scope for the present paper. What is in scope here is a different correctness question: whether the \emph{exported serving artifacts} reproduce the classifier they were exported from. Sect.~\ref{sec:parity} treats that as a release-gated, measured property, not an assumption.

\section{System}
\label{sec:arch}

\subsection{Two-Stage Pipeline with a Swappable Classifier}

Fig.~\ref{fig:architecture} shows the deployment. A coordinator host runs Kafka and replays CICIDS2017-derived flow records onto an input topic at a configurable rate. Jetson~\#1 runs the anomaly gate, a scikit-learn autoencoder and scaler trained offline (out of scope here). Flows whose reconstruction error clears a calibrated threshold are forwarded, with their anomaly score attached, to a suspicious-flow topic. Jetson~\#2 consumes that topic and classifies each flow with a 200-tree, depth-15 random forest trained in earlier work with a histogram-binning trainer and exported from its saved node table without retraining, so the board serves exactly the model selected offline. Both nodes serve a batch once it reaches a target size or once a linger timeout has elapsed since its first record, whichever comes first, so a thin stream still gets bounded latency instead of waiting to fill a batch.

\begin{figure}[t]
  \centering
  \resizebox{\textwidth}{!}{%
  \begin{tikzpicture}[
      font=\small, >=Stealth, every node/.style={align=center},
      pc/.style ={draw=gray!60,        thick, rounded corners, fill=gray!8,  minimum width=11.6cm, minimum height=1.2cm, inner sep=3pt},
      jet/.style={draw=accGreen!80!black, thick, rounded corners, fill=accGreen!8, text width=4.4cm, minimum height=1.9cm, inner sep=4pt},
      eng/.style={draw=accBlue, thick, rounded corners, fill=accBlue!5, text width=2.0cm, minimum height=0.9cm, inner sep=2pt, font=\footnotesize},
      lbl/.style={font=\footnotesize, fill=white, inner sep=1.5pt, align=center, text=accBlue!75!black},
      arr/.style ={-{Stealth[length=2.5mm]}, thick, draw=gray!60!black},
  ]
    \node[pc]  (pc)  at (0, 0.0)   {\textbf{Coordinator (Kafka)}\\[2pt]
        CICIDS2017-derived replay $\rightarrow$ \texttt{ids-network-flow}};
    \node[jet] (j1)  at (-4.2,-3.4) {\textbf{Jetson Orin Nano Super \#1}\\[2pt] anomaly gate\\[2pt]
        sklearn autoencoder + scaler\\ (joblib artifacts)};
    \node[jet] (j2)  at ( 4.2,-3.4) {\textbf{Jetson Orin Nano Super \#2}\\[2pt] classifier\\[2pt]
        feature-matrix builder $\rightarrow$\\ inference-engine contract};
    \node[eng] (onnx)  at (2.9,-5.9) {ONNX Runtime\\ (\texttt{.onnx})};
    \node[eng] (numpy) at (5.5,-5.9) {NumPy engine\\ (\texttt{.npz})};
    \draw[arr] (j2.south -| onnx)  -- (onnx.north);
    \draw[arr] (j2.south -| numpy) -- (numpy.north);
    \node[lbl] at (4.2,-6.8) {one flag: \texttt{--engine onnx\,|\,numpy}};

    \draw[arr] (pc.south -| j1) -- node[lbl]{ids-network-flow} (j1.north);
    \draw[arr] (j1.east) -- node[lbl]{ids-suspicious-flow\\ (anomaly score attached)} (j2.west);
  \end{tikzpicture}}
  \caption{Two-stage Kafka pipeline: gate on Jetson~\#1, classifier on Jetson~\#2 with two backends behind one contract, selected by a single flag. This minimal toolkit prints per-batch metrics and an optional CSV; the PostgreSQL/alerting/Grafana stack of the earlier deployment work is out of scope here.}
  \label{fig:architecture}
\end{figure}

The classifier tier is built around one contract. A feature-matrix builder turns a batch of flow records into a dense float32 matrix in a fixed column order, mapping missing and non-finite values to zero. An inference-engine interface takes that matrix and returns predictions, per-row confidences and batch timing, and has two implementations behind one command-line flag: ONNX Runtime with its CPU execution provider (GPU and TensorRT providers are selectable but not evaluated here), and a NumPy engine that loads the flattened forest of Sect.~\ref{sec:export} and walks every tree with a vectorized loop over an active-row mask, accumulating leaf class scores as rows reach leaves. Nothing else in the pipeline changes when the flag changes.

\subsection{NumPy Export from the ONNX Graph}
\label{sec:export}

The exporter treats the ONNX file as the single source of truth. It reads the \texttt{Scaler} node's offset and scale and the \texttt{TreeEnsembleClassifier} node's node table (tree and node ids, split features, thresholds, split modes, child ids) and leaf table (class ids and weights), all of which are attributes fixed by the ONNX operator specification, not by any training framework's internal format, and writes them as flat arrays to a compressed \texttt{.npz} file. Because the exporter never touches how the graph was produced, it exports a NumPy fallback for any ONNX tree ensemble that uses this operator. It supports the \texttt{BRANCH\_LEQ} split mode only and raises on any other, so an export from a framework with a different split convention fails loudly instead of producing a plausible but wrong artifact. It also refuses to export when the feature-column list embedded in the ONNX model's doc string disagrees with the externally supplied one, since a silently permuted feature order yields confident, wrong predictions instead of an obvious failure.

\subsection{Parity as a Release Gate}
\label{sec:parity}

A validation script runs both backends over the same batch-built feature matrix and reports label agreement and the maximum confidence deviation between them. It exits non-zero if agreement falls below a threshold (default: exact) or the deviation exceeds a tolerance (default $10^{-5}$), so it can gate a release instead of running as a one-off check, and it optionally writes a JSON report. On a validation run of 2{,}000 replayed CICIDS2017-derived flows the two backends agreed on every label with a maximum confidence deviation of 0.0, i.e.\ bit-identical decisions on the sampled flows. The NumPy artifact (1.23\,MB) is also $6\times$ smaller than the ONNX graph (7.38\,MB). This is the gate the software shipped with through v0.3.0, and the next section asks what it establishes. From v0.4.0 the same script also runs the boundary probes of Sect.~\ref{sec:boundary}, built from the NumPy artifact alone (Sect.~\ref{sec:discussion}).

\section{Boundary-Directed Parity Testing}
\label{sec:boundary}

Two tree runtimes fed the same row can only disagree if, at some split, they place the row on different sides of the threshold. That can only happen when the scaled feature value is within rounding distance of the threshold, and recorded traffic is almost never there. The scarcity is a property of recordings, not of the input space: evasion attacks on tree ensembles search for minimally perturbed inputs that sit exactly against the model's split boundaries~\cite{kantchelian2016evasion}, so the region replay never visits is one an adversary deliberately occupies. Replay therefore tests the runtimes where they cannot differ. We test them where they can.

For each of the forest's 95{,}741 internal nodes, with split feature $f$ and threshold $t$ in scaled space, we (i) walk from the node to its root, intersecting the constraints each ancestor split places on each feature, and write a value from inside every resulting interval into a real test-split row, so the row is routed to the node; (ii) overwrite feature $f$ with the raw float32 value that the scaler maps onto $t$, and with its float32 predecessor and successor. This yields $3 \times 95{,}741 = 287{,}223$ probes. Re-tracing each probe through the forest confirms that 95{,}664 of the 95{,}741 (99.9\%) arrive at their target node. The remainder have path intervals narrower than float32 can represent after inverse scaling. Fig.~\ref{fig:probe} sketches the two steps.

\begin{figure}[t]
  \centering
  \resizebox{\textwidth}{!}{%
  \begin{tikzpicture}[
      font=\small, >=Stealth, every node/.style={align=center},
      tn/.style ={draw=gray!60, thick, rounded corners=2pt, fill=gray!8, minimum width=1.55cm, minimum height=0.55cm, inner sep=2pt, font=\footnotesize},
      tgt/.style={tn, draw=accBlue, fill=accBlue!8, text=accBlue!75!black},
      leaf/.style={draw=gray!50, fill=white, circle, minimum size=0.32cm, inner sep=0pt},
      edge/.style={draw=gray!60, thick},
      path/.style={draw=accBlue, line width=1.4pt},
      lbl/.style={font=\footnotesize, inner sep=1.5pt, align=center},
      hdr/.style={font=\footnotesize\bfseries, align=left, anchor=west},
      probe/.style={circle, fill=accOrange, draw=white, line width=0.6pt, minimum size=0.26cm, inner sep=0pt},
      arr/.style ={-{Stealth[length=2.2mm]}, thick, draw=gray!60!black},
  ]
    % ---------------- (i) route a real row to the node ----------------
    \node[hdr] at (-0.4, 3.0) {(i) Route a real row to the node};
    \node[tn]  (r)  at (0.9, 2.3)  {$x_3 \le t_1$};
    \node[tn]  (a)  at (0.1, 1.3)  {$x_7 \le t_2$};
    \node[leaf](l1) at (1.7, 1.3)  {};
    \node[leaf](l2) at (-0.7, 0.3) {};
    \node[tgt] (n)  at (0.9, 0.3)  {$x_f \le t$};
    \node[leaf](l3) at (0.3, -0.6) {};
    \node[leaf](l4) at (1.5, -0.6) {};
    \draw[edge] (r) -- (l1);  \draw[edge] (a) -- (l2);
    \draw[edge] (n) -- (l3);  \draw[edge] (n) -- (l4);
    \draw[path] (r) -- (a);   \draw[path] (a) -- (n);
    \node[lbl, text=gray!60!black] at (-0.05, 1.95) {yes};
    \node[lbl, text=gray!60!black] at (0.85, 0.85)  {no};
    \node[lbl, anchor=west, text=accBlue!75!black] at (1.85, 0.3) {target node\\(one of 95{,}741)};
    % path constraints
    \node[lbl, anchor=north, draw=gray!40, rounded corners=2pt, fill=gray!4, inner sep=4pt, text width=4.2cm]
      (cons) at (0.9, -1.0)
      {intersect ancestor splits:\\
       $x_3 \in (-\infty, t_1]$, \quad $x_7 \in (t_2, \infty)$\\[1pt]
       write one value from each interval\\into a replay row $\Rightarrow$ row reaches $n$\\
       (99.9\% of nodes reached)};
    % ---------------- (ii) three float32 probes at the threshold ----------------
    \node[hdr] at (5.3, 3.0) {(ii) Overwrite $x_f$ with three float32 values};
    % scaled axis
    \draw[arr] (5.4, 1.4) -- (11.6, 1.4) node[lbl, anchor=west] {scaled $x_f$};
    \draw[thick, accBlue] (8.5, 1.1) -- (8.5, 1.75);
    \node[lbl, text=accBlue!75!black, anchor=south] at (8.5, 1.78) {threshold $t$};
    \node[probe] (p1) at (7.3, 1.4) {};
    \node[probe] (p2) at (8.5, 1.4) {};
    \node[probe] (p3) at (9.7, 1.4) {};
    \node[lbl, anchor=north] at (7.3, 1.2) {$t - 1$\,ulp};
    \node[lbl, anchor=north] at (8.5, 1.2) {$= t$};
    \node[lbl, anchor=north] at (9.7, 1.2) {$t + 1$\,ulp};
    \draw[gray!60, dashed] (8.5, 1.1) -- (8.5, 0.45);
    \node[lbl, anchor=south, text=gray!60!black] at (6.55, 1.55) {go left: $x_f \le t$};
    \node[lbl, anchor=south, text=gray!60!black] at (10.55, 1.55) {$x_f > t$: go right};
    % raw -> scaler -> scaled
    \node[lbl, anchor=south, text=gray!60!black] at (6.1, 2.1) {raw float32 $\xrightarrow{\;\text{scaler}\;}$ scaled};
    % float32 vs float64 comparison
    \node[lbl, anchor=north, draw=accGreen!80!black, rounded corners=2pt, fill=accGreen!8, inner sep=4pt, text width=3.1cm]
      at (6.6, 0.35)
      {\textbf{float32 scaler + compare}\\(ONNX Runtime, deployed NumPy)\\probe lands where placed:\\\textbf{0 flips} / 287{,}223};
    \node[lbl, anchor=north, draw=accOrange, rounded corners=2pt, fill=accOrange!8, inner sep=4pt, text width=3.1cm]
      at (10.5, 0.35)
      {\textbf{float64 scaler}\\scaled value differs by $\varepsilon$,\\the on-threshold probe can\\cross $t$: \textbf{328 flips}};
    \draw[arr, accOrange] (9.3, 0.35) to[bend left=10] (8.62, 1.28);
  \end{tikzpicture}}
  \caption{Boundary-directed probes. (i)~For every internal node, the constraints along its root path are intersected per feature and satisfying values are written into a real test-split row, so the row is routed to that node. (ii)~The node's split feature is then overwritten with the raw float32 values that the scaler maps onto the threshold and one float32 step either side. A traversal that scales or compares in float64 computes a slightly different scaled value and can place the on-threshold probe on the other side of~$t$.}
  \label{fig:probe}
\end{figure}

\begin{table}[htbp]
  \caption{Label disagreements with ONNX Runtime: boundary-directed probes vs.\ replaying the whole labelled test split.}
  \label{tab:boundary}
  \centering
  \small
  \begin{tabular}{l r r r}
    \toprule
    Input set & Inputs & NumPy float32 & NumPy float64 \\
    \midrule
    Probe on threshold          & 95{,}741  & 0 & 143 \\
    Probe one float32 step below & 95{,}741  & 0 & 7 \\
    Probe one float32 step above & 95{,}741  & 0 & 178 \\
    \textbf{All boundary probes} & 287{,}223 & \textbf{0} & \textbf{328} \\
    \midrule
    Unrouted per-pair probes     & 286{,}065 & 0 & 133 \\
    Replay, full test split      & 447{,}275 & 0 & 15 \\
    \bottomrule
  \end{tabular}
  \tabnote{Counts are label flips, not probability differences. Unrouted: the same boundary values written into random replay rows without the path construction, 117 rows per distinct (feature, threshold) pair. Source: \texttt{lncs/code/results/exp2\_boundary\_parity.json} and \texttt{exp2b\_probe\_ablation.json}; counts do not depend on the machine.}
\end{table}

Table~\ref{tab:boundary} compares two NumPy traversals against ONNX Runtime. They differ in one line: the dtype in which the scaler and the threshold comparison are carried out. ONNX's tree-ensemble operator holds float32 thresholds and ONNX Runtime scales and compares in float32. A traversal that does the same arithmetic in float64 computes a slightly different scaled value and can land on the other side of a threshold it was nominally equal to. The float64 variant is not a strawman: it is the natural default in NumPy, it is what an earlier version of our exporter wrote, and it passes a 2{,}000-row replay gate almost every time. At the measured replay disagreement rate of $15/447{,}275$, the chance that a 2{,}000-row replay contains no disagreeing row is $(1 - 3.35\times10^{-5})^{2000} = 93.5\%$. The boundary probes surface the same defect at a rate of $328/287{,}223$, $34\times$ denser, and deterministically: the probe set is a function of the model, not of what traffic happened to be recorded. The float32 traversal we deploy produces zero flips on all 287{,}223 probes and on all 447{,}275 replayed rows.

\textbf{Why labels flip, not just votes.} A single misrouted tree moves one vote in 200, which should almost never change a label. It does here because the trees are not independent in where they split. The trainer of the deployed forest, like every histogram-based trainer, does not search split points on raw feature values: it bins each continuous feature into at most 32 quantile bins on a sample of the training data, so at most 31 candidate thresholds per feature exist, and every tree in the forest chooses from the same candidates. The served forest shows this exactly: its 95{,}741 splits use only 815 distinct (feature, threshold) pairs, 15 of the 29 features sit at the 31-threshold cap, and a typical pair appears in 160--200 of the 200 trees. A value sitting on one such threshold sits on it in every tree that uses it, so an arithmetic disagreement there is correlated across the ensemble instead of being averaged out by it. Counting, for each probe, the trees whose leaf differs between the float32 and the float64 traversal makes this concrete: an on-threshold probe that moves any tree moves 7.5 on average, and the probes whose label flips moved 23.6. For a detector this has an operational consequence: in a fleet where some boards run the primary runtime and others the fallback, a flow whose features sit on shared thresholds, by accident or by adversarial construction, receives different verdicts depending on which board sees it. The same redundancy makes the gate cheap: 815 distinct thresholds, not 95{,}741 nodes, bound the number of boundary values that need probing.

\textbf{Is the path construction necessary?} A cheaper gate would skip step (i): write each of the 815 distinct boundary values into random replay rows and let the row's own path meet a node carrying that pair, or not. At the same probe budget this unrouted variant still finds the float64 divergence $14\times$ more densely than replay (Table~\ref{tab:boundary}), but only 38\% of its probes ever evaluate the comparison they were built to stress, it exercises 27\% of the forest's internal nodes against the construction's 99.9\%, and it finds $2.5\times$ fewer flips. Conditioned on a probe reaching a node with its pair, the two variants' flip densities agree. What the construction buys is that reaching the node happens by construction rather than by a 38\% accident, so a passing gate covers the forest's splits deterministically instead of sampling them.
\label{sec:generalize}

The probe set is derived from the artifact, so nothing ties it to this forest or this trainer. We trained four more ensembles on a 600{,}000-row sample of the same training split, with the same standardizing front end, and exported them through the standard converters (skl2onnx and onnxmltools\footnote{\url{https://github.com/onnx/sklearn-onnx}, \url{https://github.com/onnx/onnxmltools}}) into the same graph shape, a \texttt{Scaler} feeding a \texttt{TreeEnsembleClassifier}: a scikit-learn \texttt{RandomForest}~\cite{pedregosa2011sklearn} with exact midpoint splits, scikit-learn \texttt{HistGradientBoosting}, LightGBM~\cite{ke2017lightgbm} and XGBoost~\cite{chen2016xgboost} in histogram mode, whose export uses the \texttt{BRANCH\_LT} split mode. For each we compare ONNX Runtime with a float64 traversal of its graph, as above, and with the training framework's own prediction, which is the converter question a replay gate is usually written to answer. Table~\ref{tab:generalize} gives the counts.

\begin{table}[htbp]
  \caption{Boundary probes vs.\ replay on five forests trained on the same features and exported to ONNX. Flips are label disagreements with ONNX Runtime; miss is the probability that a 2{,}000-row replay gate sees no flip.}
  \label{tab:generalize}
  \centering
  \footnotesize
  \setlength{\tabcolsep}{3pt}
  \begin{tabular}{l r r r r r r r r}
    \toprule
     & & & \multicolumn{3}{c}{float64 traversal} & \multicolumn{3}{c}{training framework} \\
    \cmidrule(lr){4-6}\cmidrule(lr){7-9}
    Forest (bins) & Pairs & Trees/pair & probes & replay & miss & probes & replay & miss \\
    \midrule
    deployed RF (32)       & 815      & 72.1 & 328 & 15      & 93.5\% &     &         &        \\
    sklearn RF (exact)     & 57{,}373 & 2.1  & 182 & 8       & 96.5\% & 64  & 4       & 98.2\% \\
    sklearn HGB (255)      & 1{,}370  & 4.3  & 233 & 895     & 1.8\%  & 195 & 1{,}362 & 0.2\%  \\
    LightGBM (255)                & 1{,}538  & 3.8  & 90  & 31      & 87.1\% & 75  & 11      & 95.2\% \\
    XGBoost hist (256)            & 2{,}224  & 3.9  & 429 & 1{,}066 & 0.8\%  & 132 & 687     & 4.6\%  \\
    \bottomrule
  \end{tabular}
  \tabnote{sklearn: scikit-learn. Pairs: distinct (feature, threshold) pairs; trees/pair: trees sharing a pair, mean. Probes: three per internal node, over all 95{,}741 / a 100{,}000-node sample of 128{,}315 / all 6{,}000 / 6{,}000 / 9{,}070 nodes; 99.9--100\% reach their node. Replay: the 447{,}275-row test split. The float32 traversal flips no label on any probe or replayed row of any forest. Source: \texttt{lncs/code/results/exp5\_generalization.json}.}
\end{table}

The float32 traversal agrees with ONNX Runtime on every probe and every replayed row of every forest. For the three forests whose divergence is rare, the two random forests and LightGBM, a 2{,}000-row replay passes 87 to 98\% of the time, and the probes find the same divergence 24 to $169\times$ more densely. For HistGradientBoosting and XGBoost the divergence is common: 0.2 to 0.3\% of replayed rows change label and a tree moves on 50 to 70\% of them, so replay catches it and the probes add only 3 to $7\times$. For those two exports exactness is already lost, and the open question is whether a disagreement rate of that size is acceptable. The mechanism of Sect.~\ref{sec:boundary} also needs one qualification. How many trees share a threshold depends on how many bins the trainer uses: the deployed forest's 32 bins give 72 trees per pair, the 255-bin trainers about 4, and exact splits 2. Labels still flip in every forest, because a flip needs a boundary value to move many trees, and shared thresholds are only one way to get there. In the scikit-learn forest an on-threshold probe that moves any tree moves 2.2 on average, yet the probes whose label flipped had moved 15. The extra moves come from the base row's other features: real traffic places them within rounding distance of some threshold often enough that 0.36\% of unmodified replay rows move a tree too. Sharing does raise the chance that a moved tree flips the label, 0.29\% against 0.16\%, and with 815 pairs against 57{,}373 it is what makes the exhaustive gate cheap.

\textbf{What replay is for.} ONNX Runtime 1.26.0 and 1.27.0 label the scikit-learn export differently from the five other releases we tried on it (1.19.2, 1.22.1, 1.25.1, 1.28.0 and 1.30.0). The graph carries $P(\text{attack})$ as a single weight column, and those two releases report attack whenever it is positive instead of above one half, which relabels 380{,}842 of the 447{,}275 test flows (85\%) and returns a negative first probability column. The regression is known upstream: a fix for a LOGISTIC post-transform bug removed the runtime's all-positive-weights flag in 1.26.0, the RandomForest breakage was reported two and a half months later, and the flag came back in 1.28.0\footnote{onnxruntime pull requests 27536 and 28685 and issue 28557; the ONNX reference evaluator labels the same graph as scikit-learn does.}. A replay of any size finds this immediately. The boundary probes cannot, because they compare traversals of one graph under one runtime, and the deployed forest escapes it only because its export carries both class columns. A parity gate therefore needs both: replay for gross changes in the runtime, probes for the divergences replay cannot see.

\section{Experimental Setup}

The testbed is two NVIDIA Jetson Orin Nano Super Developer Kits (6-core Arm Cortex-A78AE CPU, Ampere GPU, 8\,GB LPDDR5) and a coordinator host running Kafka, all three on one Wi-Fi LAN (median round trip from the coordinator to either board 5--8\,ms during the runs). The link affects only the end-to-end figures of Sect.~\ref{sec:e2e}, not the in-process measurements elsewhere in Sect.~\ref{sec:results}. \texttt{scripts/csv\_producer.py} replays CICIDS2017-derived~\cite{sharafaldin2018toward} flow rows onto the input topic at a configurable rate. \texttt{gate\_node.py} runs on Jetson~\#1 and \texttt{classifier\_node.py} on Jetson~\#2, the latter re-run once per backend (\texttt{--engine onnx} and \texttt{--engine numpy}) with the gate held fixed, and once more across a batch-size sweep (1, 20, 500) to characterize batch sensitivity. Reported metrics: throughput (rows/s), p95 per-row latency (ms), resident memory (RSS), board power (W, \texttt{tegrastats}, not idle-subtracted, so figures are total-board, not active-only), cold-start time to first verdict, and the parity metrics of Sect.~\ref{sec:parity}.

\textbf{Off-board experiments.} The boundary-parity and shift experiments (Tables~\ref{tab:boundary},~\ref{tab:drift}) run the same served artifacts, with the code under \texttt{lncs/code/}, on an ARM64 development machine (Apple M3, NumPy 2.5, ONNX Runtime 1.30, CPU provider), not on the boards. The traversal and early-exit experiments (Fig.~\ref{fig:lockstep}, Table~\ref{tab:earlyexit}) ran there first and were then re-run on Jetson~\#2 with the same code and test split (25\,W power mode, NumPy 2.2.6, ONNX Runtime 1.23.2, the versions the board serves with). Their label counts match the M3 run exactly, and the tables report the board. The forests of Table~\ref{tab:generalize} were trained there too (scikit-learn 1.9, LightGBM 4.7, XGBoost 3.4, skl2onnx 1.20, onnxmltools 1.16; 200 trees each, depth 15 for the random forest and 8 for XGBoost, library defaults otherwise) and the version sweep installed each ONNX Runtime release into its own environment. They use the labelled CICIDS2017 test split (447{,}275 flows, 15.0\% attacks) and, for the shift experiment, the CSE-CIC-IDS2018 test split prepared identically (477{,}771 flows, 12.7\% attacks). Disagreement and label-identity counts do not depend on the machine. Throughput does, and is labelled as such wherever it appears.

\section{Results}
\label{sec:results}

\subsection{Steady-State Cost per Backend}

Table~\ref{tab:engines} reports classifier-node cost at batch size 20 for each backend, with the gate's own cost reported separately since it runs a different model (an autoencoder, not the forest). Over five runs per backend, ONNX Runtime sustains roughly $119\times$ the throughput of the per-tree NumPy backend at this batch size and a fraction of its p95 latency, at the cost of $3.3\times$ the resident memory and at marginally lower measured board power (4.66 vs.\ 4.99\,W; since power is not idle-subtracted, this difference should be read cautiously rather than as a per-inference energy claim). The per-tree row is the traversal of the released version (v0.2.0); the block-lockstep row is the traversal of Sect.~\ref{sec:lockstep} served by the same node, $22\times$ the per-tree row and about a fifth of ONNX Runtime, at the per-tree row's memory.

\begin{table}[htbp]
  \caption{Classifier-node steady-state cost per backend (batch size 20, mean $\pm$ SD over five runs on Jetson~\#2) and gate-node cost (batch size 100, one run on Jetson~\#2).}
  \label{tab:engines}
  \centering
  \small
  \begin{tabular}{l c c c c}
    \toprule
    Node / engine & Throughput (rows/s) & p95 (ms/row) & RSS (MB) & Power (W) \\
    \midrule
    Classifier (ONNX Runtime) & 15{,}174 $\pm$ 179 & 0.0847 $\pm$ 0.0024 & 156.8 $\pm$ 1.0 & 4.66 \\
    Classifier (NumPy, per-tree) & 127 $\pm$ 1 & 9.50 $\pm$ 0.08 & 47.6 $\pm$ 0.9 & 4.99 \\
    Classifier (NumPy, lockstep) & 2{,}752 $\pm$ 17 & 0.412 $\pm$ 0.005 & 47.3 $\pm$ 0.1 & 4.67 \\
    Gate (sklearn autoencoder) & 21{,}191 & 0.0569 & 123.6 & 5.25 \\
    \bottomrule
  \end{tabular}
  \tabnote{Power is total-board \texttt{tegrastats} power (\texttt{VDD\_IN}), not idle-subtracted. Gate row: \texttt{gate\_node.py} at its default batch size serving a 20{,}000-record Kafka replay; throughput and p95 are its own inference time, RSS is the node process's, and the board drew 4.75\,W idle before the replay (\texttt{lncs/code/results\_jetson/exp6\_gate\_node.json}). Its forward rate on that replay was 11.8\%, close to the test split's 11.3\% (Sect.~\ref{sec:drift}). The software's earlier validation run recorded 13{,}915\,rows/s and 0.0783\,ms/row for the gate on Jetson~\#1. Classifier rows: 30{,}000-record replays at about 550\,rows/s, 25\,W power mode, throughput from the node's own inference timer (\texttt{lncs/code/results\_jetson/exp8\_summary.json}); power SDs are at most 0.02\,W. The software's first single validation run had recorded 20{,}185\,rows/s, p95 0.0676\,ms and 155.6\,MB for ONNX Runtime and 152\,rows/s, p95 8.2267\,ms and 46.2\,MB for per-tree NumPy; the repeats do not reproduce its ONNX throughput, and its power mode was not recorded.}
\end{table}

\subsection{Cold Start: Time to First Verdict}

Table~\ref{tab:coldstart} breaks down the time from process start to the first served verdict. ONNX Runtime pays a larger one-time session-construction cost (373\,ms) than NumPy's artifact load (80\,ms), but performs its own first prediction fast (7\,ms), whereas the per-tree NumPy traversal needs 114\,ms for one row because a single row cannot amortize its per-tree dispatch. NumPy still reaches its first verdict sooner overall (482 vs.\ 670\,ms), since its faster load dominates over the whole warm-up path. The block-lockstep traversal removes the first-call penalty (2.2\,ms) and brings the first verdict to 357\,ms. This favors NumPy specifically for restart-sensitive scenarios (e.g.\ frequent container restarts or fast failover) even though ONNX Runtime is far ahead on steady-state throughput.

\begin{table}[htbp]
  \caption{Cold-start latency per backend, page cache dropped before every run: mean $\pm$ SD of five runs on each board (lockstep: five runs on Jetson~\#2). Time to first verdict runs from launching the process.}
  \label{tab:coldstart}
  \centering
  \small
  \begin{tabular}{l c c c}
    \toprule
    Engine & Load/session build (ms) & First prediction (ms) & Time to first verdict (ms) \\
    \midrule
    ONNX Runtime        & 373 $\pm$ 4 & 7.0 $\pm$ 1.7   & 670 $\pm$ 23 \\
    NumPy, per-tree     & 80 $\pm$ 2  & 113.9 $\pm$ 1.8 & 482 $\pm$ 17 \\
    NumPy, lockstep     & 79 $\pm$ 0  & 2.2 $\pm$ 0.1   & 357 $\pm$ 1 \\
    \bottomrule
  \end{tabular}
\end{table}

\subsection{End-to-End Latency over Wi-Fi}
\label{sec:e2e}

Table~\ref{tab:e2e} follows every flow the gate forwards from the producer's send timestamp to the classifier's verdict, with the gate on Jetson~\#1 and the ONNX classifier on Jetson~\#2, both at their default batch size and a 1\,s linger. Each board's timestamps are moved onto the coordinator's clock by an offset measured over SSH before and after the run (error and drift both under 2\,ms). At a light load, latency is set by the batching policy rather than by inference or the network: the classifier receives about 7 forwarded flows per second, never fills a batch of 20, and waits out its full linger on almost every flow, so its stage alone takes 1.06\,s at the median. At the producer's maximum rate the gate drains about 965\,rows/s, far below the 21{,}191\,rows/s its inference sustains (Table~\ref{tab:engines}), so flows queue ahead of it and its stage dominates, while the classifier stage falls to 0.16\,s because forwarded flows now fill its batches. A deployment that needs sub-second verdicts at low load should shorten the classifier's linger before touching the model.

\begin{table}[htbp]
  \caption{Producer-to-verdict latency of gate-forwarded flows over Wi-Fi (gate on Jetson~\#1, ONNX classifier on Jetson~\#2), p50\,/\,p95 in seconds.}
  \label{tab:e2e}
  \centering
  \small
  \setlength{\tabcolsep}{7pt}
  \begin{tabular}{l r c c c}
    \toprule
    Offered load & Verdicts & Gate stage & Classifier stage & End-to-end \\
    \midrule
    64\,rows/s, 6{,}000 records         & 685     & 0.56\,/\,1.03 & 1.06\,/\,1.16 & 1.58\,/\,2.11 \\
    $\approx$1{,}470\,rows/s, 20{,}000 records & 2{,}362 & 3.38\,/\,6.31 & 0.16\,/\,0.29 & 3.55\,/\,6.44 \\
    \bottomrule
  \end{tabular}
  \tabnote{One run per load; the producer was asked for 100\,rows/s and for its maximum rate. Every forwarded flow (11.4\% and 11.8\% of the records) received a verdict. Gate stage: send to forward; classifier stage: forward to verdict. Source: \texttt{lncs/code/results\_jetson/exp7\_e2e.json}.}
\end{table}

\subsection{NumPy Batch-Size Sensitivity}
\label{sec:batch-sensitivity}

The original NumPy engine traverses each tree with its own Python-level \texttt{while}-loop, one tree at a time. Because each such loop's per-call NumPy dispatch overhead is paid once per tree regardless of how little work that call does, its throughput scales strongly with batch size on Jetson hardware: 9\,rows/s at batch 1, 150\,rows/s at batch 20, and 1{,}650\,rows/s at batch 500 on the same board (\texttt{submission\_checklist.md}, this software's own validation run). ONNX Runtime is comparatively insensitive to batch size over the same range (Sect.~\ref{sec:discussion}). This means the two backends should not be compared at a single fixed batch size without stating it, and a deployment choosing the NumPy fallback should size its batch/linger policy accordingly rather than default to a small batch tuned for the ONNX path. Or, as Sect.~\ref{sec:lockstep} shows, the per-tree loop itself is the thing to fix at those batch sizes.

\subsection{Block-Lockstep Traversal}
\label{sec:lockstep}

The batch-size sensitivity in Sect.~\ref{sec:batch-sensitivity} is a symptom of the per-tree Python loop, not of a NumPy-only backend. Node ids are globally unique across the flattened forest, so nothing requires visiting trees one after another: a (trees $\times$ rows) matrix of current node ids can be advanced one level per iteration, for every tree and every row at once, cutting the Python-level loop from $O(\text{trees} \times \text{depth})$ iterations to $O(\text{depth})$.

Our first version did exactly that over all 200 trees, and we benchmarked it on synthetic rows. Re-measured on real test-split flows it was 37\% \emph{slower} than the per-tree loop at batch size 2{,}000: the $200 \times 2{,}000$ node matrix no longer fits in cache, and every level becomes a scattered gather over it. The version reported here advances blocks of about 8{,}192 (tree, row) pairs at a time, so the working set stays cache-sized at any batch size, and accumulates leaf scores with a gather-and-sum rather than a scatter-add. Fig.~\ref{fig:lockstep} gives throughput on real flows, best of five runs. It times the traversal alone, so its per-tree series runs about twice the serving-node figures of Sect.~\ref{sec:batch-sensitivity}.

\begin{figure}[t]
  \centering
  \begin{tikzpicture}[font=\small]
    \begin{groupplot}[
        group style={group size=2 by 1, horizontal sep=0.9cm, ylabels at=edge left, yticklabels at=edge left},
        width=0.505\textwidth, height=5.2cm,
        xmode=log, ymode=log, log basis x=10, log basis y=10,
        xmin=0.6, xmax=8000, ymin=8, ymax=60000,
        clip=false,
        xtick={1,20,100,500,2000}, xticklabels={1,20,100,500,2{,}000},
        ytick={10,100,1000,10000}, yticklabels={10,100,1{,}000,10{,}000},
        minor tick num=0, log ticks with fixed point,
        xlabel={Batch size (rows)}, ylabel={Throughput (rows/s)},
        axis line style={gray!60}, tick style={gray!60}, tick label style={font=\footnotesize, text=gray!40!black},
        label style={font=\footnotesize, text=gray!40!black},
        grid=major, grid style={gray!18, dashed},
        title style={font=\footnotesize\bfseries, yshift=-2pt},
        every axis plot/.append style={line width=1.2pt, mark size=2.2pt},
        legend style={at={(0.97,0.03)}, anchor=south east, font=\footnotesize, draw=gray!40, fill=white, fill opacity=0.9, text opacity=1, cells={anchor=west}, row sep=-1pt},
        legend cell align=left,
      ]
      % ---- Jetson #2 ----
      \nextgroupplot[title={Jetson Orin Nano Super (\#2), 25\,W}]
        \addplot[color=accBlue, mark=*, mark options={fill=accBlue, draw=white, line width=0.5pt}]
          coordinates {(1,1435.3) (20,6325.8) (100,6594.4) (500,7112.9) (2000,7059.7)};
        \addlegendentry{block-lockstep}
        \addplot[color=accOrange, mark=square*, mark options={fill=accOrange, draw=white, line width=0.5pt}]
          coordinates {(1,17.8) (20,324.7) (100,1083.9) (500,3292.7) (2000,6020.1)};
        \addlegendentry{per-tree loop}
        \node[font=\scriptsize, text=accBlue!75!black, anchor=north west, xshift=-2pt, yshift=-2pt] at (axis cs:1,1435.3) {$80.5\times$};
        \node[font=\scriptsize, text=accBlue!75!black, anchor=south] at (axis cs:20,6325.8)   {$19.5\times$};
        \node[font=\scriptsize, text=accBlue!75!black, anchor=south] at (axis cs:100,6594.4)  {$6.1\times$};
        \node[font=\scriptsize, text=accBlue!75!black, anchor=south] at (axis cs:500,7112.9)  {$2.2\times$};
        \node[font=\scriptsize, text=accBlue!75!black, anchor=west, xshift=3pt] at (axis cs:2000,7059.7) {$1.2\times$};
      % ---- Apple M3 ----
      \nextgroupplot[title={Apple M3 (development machine)}]
        \addplot[color=accBlue, mark=*, mark options={fill=accBlue, draw=white, line width=0.5pt}]
          coordinates {(1,5539.9) (20,15807.1) (100,16051.0) (500,16141.7) (2000,15849.5)};
        \addplot[color=accOrange, mark=square*, mark options={fill=accOrange, draw=white, line width=0.5pt}]
          coordinates {(1,78.1) (20,1424.0) (100,4509.1) (500,11770.9) (2000,18176.8)};
        \node[font=\scriptsize, text=accBlue!75!black, anchor=north west, xshift=-2pt, yshift=-2pt] at (axis cs:1,5539.9) {$71\times$};
        \node[font=\scriptsize, text=accBlue!75!black, anchor=south] at (axis cs:20,15807.1)  {$11.1\times$};
        \node[font=\scriptsize, text=accBlue!75!black, anchor=south] at (axis cs:100,16051.0) {$3.6\times$};
        \node[font=\scriptsize, text=accBlue!75!black, anchor=south] at (axis cs:500,16141.7) {$1.4\times$};
        \node[font=\scriptsize, text=accOrange!80!black, anchor=west, xshift=3pt] at (axis cs:2000,18176.8) {$0.9\times$};
    \end{groupplot}
  \end{tikzpicture}
  \caption{Per-tree loop vs.\ block-lockstep traversal of the same artifact on real test-split flows, best of five runs, log--log axes shared by both panels. Labels give the lockstep speedup at each batch size. On Jetson~\#2 lockstep traverses 1{,}435 rows/s at batch~1 and 7{,}060 at 2{,}000 against 18 and 6{,}020 for the per-tree loop; on the M3 the per-tree loop overtakes it at 2{,}000 (18{,}177 vs.\ 15{,}850 rows/s). Labels match the per-tree loop on 5{,}000/5{,}000 rows and ONNX Runtime on 447{,}275/447{,}275. Source: \texttt{lncs/code/results\_jetson/exp1\_lockstep\_benchmark.json} and \texttt{lncs/code/results/}.}
  \label{fig:lockstep}
\end{figure}

Block-lockstep throughput is nearly flat in batch size, which is the property a streaming node wants: its batches are set by a linger timeout, not chosen for the engine's convenience (the classifier node defaults to 20, the gate to 100). On the board, lockstep stays ahead at every batch size measured, though its lead shrinks to $1.2\times$ at 2{,}000, where the per-tree loop's dispatch overhead is already amortized and its one-tree working set is the most cache-friendly of all. On the M3 the per-tree loop overtakes it there ($0.9\times$), so the crossover depends on the host, and a deployment that serves batches of that size should measure both on its own hardware. The rewrite changes no label: it agrees with ONNX Runtime on every row of the test split, and the software's release gate passes with a maximum confidence difference of exactly zero.

\subsection{Exact Early Exit on Gate-Forwarded Traffic}
\label{sec:earlyexit}

For the binary forest let $m$ be the running margin (attack score minus benign score) after some trees have been evaluated, and let $[\ell_t, u_t]$ be the range of tree $t$'s leaf margins, read once from the artifact. If $m + \sum_{t\,\text{unevaluated}} \ell_t > 0$ the label is attack whatever the remaining trees say. If $m + \sum u_t \le 0$ it is benign. Rows meeting either condition leave the batch. The rest continue with the next block of trees. Labels are identical to the full forest's by construction. With equally weighted trees no row can qualify before more than half have voted, so the saving is capped at $2\times$. We check first after 104 trees and then every 8.

\begin{table}[htbp]
  \caption{Exact early exit vs.\ full forest, batch size 2{,}000, on Jetson~\#2.}
  \label{tab:earlyexit}
  \centering
  \small
  \begin{tabular}{l r r r r}
    \toprule
    Traffic & Rows & Labels identical & Mean trees used & Speedup \\
    \midrule
    Forwarded by the gate & 50{,}548  & 50{,}548  & 104.1 & $1.9\times$ \\
    All test-split flows  & 447{,}275 & 447{,}275 & 105.0 & $1.9\times$ \\
    \bottomrule
  \end{tabular}
  \tabnote{Source: \texttt{lncs/code/results\_jetson/exp3\_exact\_early\_exit.json}. 99.3\% of forwarded rows and 95.2\% of all rows stop at the first check. On the Apple M3 the speedups are $2.05\times$ and $1.92\times$; every other column is identical.}
\end{table}

Table~\ref{tab:earlyexit} shows the rule running at its ceiling: 99.3\% of the flows the gate forwards are settled at the first opportunity. The forest is near-unanimous on the traffic it serves, and what limits the saving is the exactness requirement, not the data. Statistical stopping rules~\cite{hernandezlobato2009pruning,wang2018qwyc} would stop far earlier on the same flows, at the price of verdicts that can differ between the primary and the fallback runtime, the property Sect.~\ref{sec:boundary} exists to protect.

\subsection{What Survives Distribution Shift}
\label{sec:drift}

Nothing in this paper adapts the detector to drift~\cite{gama2014survey}, so it is fair to ask which of its results still hold once traffic stops resembling the training data. We freeze the deployed forest, gate and threshold (all fitted on CICIDS2017) and run them unchanged on the CSE-CIC-IDS2018 test split: 477{,}771 flows from a different network, year and attack toolset, carrying the same 29 features. This is a cross-dataset shift, harsher than gradual drift, and we use it as a stress test rather than a drift model.

\begin{table}[htbp]
  \caption{The deployed model and gate, frozen, on the in-distribution (CICIDS2017) and shifted (CSE-CIC-IDS2018) test splits.}
  \label{tab:drift}
  \centering
  \small
  \begin{tabular}{l r r}
    \toprule
     & CICIDS2017 & CSE-CIC-IDS2018 \\
    \midrule
    Flows                                         & 447{,}275 & 477{,}771 \\
    \multicolumn{3}{l}{\emph{Detection quality (moves with the data)}} \\
    \quad Classifier attack F1 / recall           & 0.983 / 0.993 & 0.238 / 0.193 \\
    \quad Gate attack recall                      & 0.723 & 0.009 \\
    \quad Gate forward rate (benign only)         & 0.113 (0.005) & 0.025 (0.027) \\
    \quad Two-stage end-to-end attack F1          & 0.839 & 0.011 \\
    \multicolumn{3}{l}{\emph{Serving guarantees (do not move)}} \\
    \quad float32 vs.\ ORT, replay flips  & 0 & 0 \\
    \quad float32 vs.\ ORT, on-threshold probe flips & 0 / 95{,}741 & 0 / 95{,}741 \\
    \quad float64 vs.\ ORT, replay / probe flips & 15 / 143 & 6 / 125 \\
    \quad Early-exit labels identical to full forest & 447{,}275 & 477{,}771 \\
    \multicolumn{3}{l}{\emph{Early-exit statistics (label-free)}} \\
    \quad Mean trees used, all / forwarded        & 105.0 / 104.1 & 108.2 / 132.2 \\
    \quad Undecided at first check, all / forwarded & 4.8\% / 0.7\% & 28.2\% / 68.1\% \\
    \bottomrule
  \end{tabular}
  \tabnote{Source: \texttt{lncs/code/results/exp4\_drift.json}. float32/float64: dtype of the NumPy traversal; ORT: ONNX Runtime. Probe rows take their unconstrained features from the dataset in that column; the probed thresholds are the model's and do not change.}
\end{table}

Table~\ref{tab:drift} separates three kinds of result. \emph{Detection quality collapses}: classifier F1 falls from 0.983 to 0.238, and the anomaly gate, whose threshold is a quantile of 2017 benign traffic, passes under 1\% of 2018 attacks while forwarding five times more benign flows than before, so the two-stage pipeline detects almost nothing end to end. Several 2018 attack families (HOIC, Hulk, SSH brute force) are missed by the classifier entirely. This agrees with the cross-dataset failures reported for these datasets and is not something a serving layer can repair. \emph{The serving guarantees are untouched}: the float32 traversal agrees with ONNX Runtime on every shifted flow and on every threshold probe built from shifted rows, and exact early exit returns the full forest's label on all 477{,}771 flows. They are properties of the model and the arithmetic, not of the traffic. The boundary gate has a further advantage here: a replay gate is only as current as its recording, whereas the probe set is regenerated from each retrained model's split table and needs no traffic at all. Note also that replay finds even fewer float64 flips on shifted traffic (6) than in distribution (15). Replay-based parity evidence gets weaker exactly when the model is about to be retrained.

\emph{The early-exit statistics move, without labels.} The share of flows the forest cannot settle at the first check rises from 4.8\% to 28.2\% (from 0.7\% to 68.1\% on forwarded traffic): under shift the trees stop agreeing with each other. This is the margin-density signal that Sethi and Kantardzic~\cite{sethi2017md3} propose for unlabeled drift detection, and here it is a by-product the serving path already computes. Table~\ref{tab:alarm} turns it into an alarm. We take the undecided rate over windows of $W$ consecutive flows, set the threshold to the largest rate seen on the first half of the in-distribution windows, and count false alarms on the other half and detections on the shifted windows. At $W=10{,}000$ the threshold is 6.9\%, no held-out window reaches it and every shifted window does (the lowest at 18.3\%). At $W=2{,}000$ the in-distribution windows spread to 12.1\% and only 140 of 238 shifted windows trip the alarm, because the 2018 split contains segments the forest settles as confidently as 2017 traffic. On forwarded flows, which is what the classifier node sees, 500 forwarded flows suffice: threshold 3.4\%, no false alarm on 51 held-out windows, 23 of 23 shifted windows detected. The windows are consecutive rows of the splits, whose order follows the captures' label phases (2--22\% attacks per 10{,}000-row window in 2017), so they are traffic segments with their own attack mix rather than timed traffic, and the calibration covers one shift. The alarm also cannot see drift that leaves the forest confidently wrong, as it is for the missed attack families above.

\begin{table}[htbp]
  \caption{The undecided rate as a drift alarm: threshold = largest rate on the calibration half of the in-distribution windows.}
  \label{tab:alarm}
  \centering
  \footnotesize
  \begin{tabular}{l r r r r r}
    \toprule
    Traffic & $W$ (flows) & In-dist.\ min--max & Threshold & False alarms, held out & Detections, shifted \\
    \midrule
    All      & 2{,}000  & 0.0--12.1\%  & 12.1\% & 0 / 112 & 140 / 238 \\
    All      & 10{,}000 & 3.1--6.9\%   & 6.9\%  & 0 / 22  & 47 / 47 \\
    Forwarded & 100     & 0.0--10.0\%  & 10.0\% & 0 / 253 & 107 / 119 \\
    Forwarded & 500     & 0.0--3.4\%   & 3.4\%  & 0 / 51  & 23 / 23 \\
    \bottomrule
  \end{tabular}
  \tabnote{Source: \texttt{lncs/code/results/exp9\_drift\_alarm.json}. In-distribution min--max is over all windows of that size; the threshold is the maximum over the calibration half only. Shifted windows at $W=10{,}000$ range from 18.3\% to 37.3\%, forwarded at $W=500$ from 50.0\% to 91.8\%.}
\end{table}

\section{Discussion}
\label{sec:discussion}

\textbf{Inference is no longer the bottleneck.} Once the ONNX Runtime backend is in use, classifier-node throughput (15{,}174\,rows/s) is far above what the pipeline delivers: at the producer's maximum rate the gate drains about 965\,rows/s (Sect.~\ref{sec:e2e}) although its own inference scores 21{,}191\,rows/s (Table~\ref{tab:engines}), so the ceiling sits in the node's consume-and-forward loop, which runs the pure-Python \texttt{kafka-python} client. We have not isolated the client from the rest of that loop. A C-backed client such as \texttt{confluent-kafka} is the next lever to test, and an end-to-end throughput figure should not be attributed to the classifier without that control.

\textbf{ONNX vs.\ NumPy is a trade-off, not a strict ordering.} Table~\ref{tab:engines} makes ONNX Runtime look like the only reasonable choice, but Table~\ref{tab:coldstart} and the artifact sizes of Sect.~\ref{sec:parity} complicate that: the NumPy artifact is $6\times$ smaller and reaches its first verdict sooner, which matters on boards where \texttt{onnxruntime} cannot be installed at all, or where restart frequency dominates steady-state throughput. The toolkit's role is to make that trade-off a one-flag choice rather than a rewrite.

\textbf{What replay-based parity gates miss.} The software's release gate replays 2{,}000 rows and had never failed. Sect.~\ref{sec:boundary} shows that this says little: a traversal that disagrees with ONNX Runtime at split boundaries passes such a gate 93.5\% of the time. The practical change is small, and v0.4.0 of the software makes it: \texttt{validate\_parity.py} derives the probe set from the artifact's own split table, holding the features a path does not constrain at the scaler offset so that no dataset is needed, and runs it alongside the replay. On the shipped forest that is 287{,}223 probes, 99.9\% of which reach their node, with 0 label disagreements and a maximum confidence deviation of $4\times10^{-7}$ between the two backends, in 16\,s on the development machine. A float64 fallback fails the same gate. Because the forest has only 815 distinct thresholds, an exhaustive probe of every distinct boundary value is cheaper than the replay it complements. Sect.~\ref{sec:generalize} bounds how far this carries: the probes find divergence on every forest we exported, but the exhaustive gate is cheap only where thresholds are shared, which is a matter of bin count (815 pairs at 32 bins, 57{,}373 for exact splits), and where a forest already disagrees with its export on 0.3\% of traffic, replay finds that without help. The version regression of ONNX Runtime 1.26 and 1.27 is the other half of the argument: probes do not replace replay, they cover what it cannot.

\textbf{Lockstep traversal and the lesson from ARM tree-ensemble work.} The speedups of QuickScorer~\cite{lucchese2015quickscorer} and of its ARM port by Koschel et al.~\cite{koschel2023arm} come from evaluating many trees per instruction with bitvectors. We took only the underlying lesson, stop walking trees one at a time, and applied it at the NumPy level. Fig.~\ref{fig:lockstep} is the result, including the first attempt's failure: whole-forest lockstep looked good on synthetic rows and lost to the old loop on real ones at large batches, until the traversal was blocked to cache size. NEON intrinsics and fixed-point thresholds remain a separate, larger step. InTreeger~\cite{bart2025intreeger} takes that step furthest: it emits integer-only C for a tree ensemble and claims no loss of precision. The probes of Sect.~\ref{sec:boundary} are the direct way to check such a claim, because integer thresholds move the boundary question from rounding to the fixed-point conversion. ONNX-level quantization of the tree graph is a similar open step. Eren et al.~\cite{eren2026drdos} show that it matters for network-security classifiers on deep architectures, but it has not been evaluated on a \texttt{TreeEnsembleClassifier}, whose cost is branching and memory access rather than matrix multiplies.

\textbf{Limitations.} The classifier and cold-start figures (Tables~\ref{tab:engines},~\ref{tab:coldstart}) are means over five runs in the 25\,W power mode, and the boards' power mode changes absolute throughput. The first validation run's ONNX throughput (20{,}185\,rows/s) was not reproduced. The gate row, the end-to-end runs and the replay parity check of Sect.~\ref{sec:parity} are single runs. The parity counts are deterministic, the others are not. Disagreement and label-identity counts (Tables~\ref{tab:boundary},~\ref{tab:generalize},~\ref{tab:earlyexit}) are machine-independent for the NumPy engine but were taken against one ONNX Runtime build (1.30, CPU provider), except the release sweep of Sect.~\ref{sec:generalize}, which covers one export only. A different provider is a different runtime and needs its own probe run. Power is total-board \texttt{tegrastats} power without an idle baseline, so it is not per-inference energy. End-to-end latency was measured over Wi-Fi only, one run per load. A wired link would shrink the transport share, which is already a few milliseconds next to the batching delays. The probes cover axis-aligned splits of five forests, \texttt{BRANCH\_LEQ} and \texttt{BRANCH\_LT}, but not missing-value routing or categorical splits, and the shipped NumPy exporter still accepts \texttt{BRANCH\_LEQ} only. The traversal that handles every split mode lives in the experiment code. Finally, the detector itself is static. On its own test split the deployed gate forwards only 72.3\% of attacks, which bounds end-to-end recall however the classifier is served, and under the shift of Sect.~\ref{sec:drift} both gate and classifier fail. This paper measures that failure but offers no adaptation, and the undecided-rate alarm is calibrated on a single shift, on windows of consecutive split rows rather than timed traffic.

\section{Conclusion}

A detector that can be served by two runtimes is only as trustworthy as the evidence that they agree, and a 2{,}000-row replay is weak evidence: a float64 traversal that flips 328 labels on boundary-directed probes flips 15 on the entire 447{,}275-row test split and passes that gate 93.5\% of the time. Probing every split threshold of the deployed forest gives a deterministic answer instead, and for the float32 traversal we deploy the answer is zero disagreements. The same probes find float64 and converter divergence on four other forests 3 to $169\times$ more densely than replay, while replay is what exposes two ONNX Runtime releases that relabel 85\% of flows for a stock scikit-learn export. A parity gate needs both. Around that result, a block-lockstep rewrite makes the NumPy fallback usable at streaming batch sizes and an exact early-exit rule halves its work without being able to change a verdict. None of this makes the detector robust to drift: on CSE-CIC-IDS2018 the frozen model and gate collapse, while parity and early-exit exactness hold on every shifted flow and the early-exit undecided rate rises sixfold without labels. Beyond that: the undecided-rate alarm on gradual drift and timed traffic, boundary probes for other ONNX Runtime providers and for missing-value and categorical splits, NEON-level traversal~\cite{koschel2023arm}, ONNX-level quantization of the tree graph~\cite{eren2026drdos}, and a C-backed Kafka client.

\begin{credits}
\subsubsection{\discintname}
The author declares that he has no competing interests.
\end{credits}

\bibliographystyle{splncs04}
\bibliography{references}

\end{document}
